% ---------------------------------------------------------------------------
%  05-storage.tex -- the storage layer: why it is a parallelism story
% ---------------------------------------------------------------------------
\section{Storage: memory bandwidth is a shared resource}
\label{sec:storage}

It is tempting to treat compression as an orthogonal memory optimisation.  It
is not, and this is one of the framing points worth making explicitly in a
paper: on a many-core machine the last-level cache and the memory controller
are \emph{shared}, so the resident size of the passed set is a parallelism
parameter.  A search whose working set spills to DRAM does not scale, however
perfect its locking is, because every worker is then queued behind the same
memory bus.

Three mechanisms shrink the resident set.

\subsection{Zones stored as minimal graphs}

\FigMingraph

A canonical DBM is a shortest-path closure, so most of its entries are implied
by the others.  The \emph{minimal graph} keeps only the non-redundant
constraints, with $16$-bit bounds when they fit: for a $21$-clock zone that is
$\approx 144$ bytes instead of $1764$, roughly $12\times$.  The codec
(\file{HDBM/include/ptctlmc/mindbm.hpp}) is a port of \uppaal{}'s UDBM
\code{mingraph}, re-expressed in \tool{}'s raw bound encoding and written as
\code{\_\_host\_\_ \_\_device\_\_} functions so the GPU backends compile the same
codec the CPU unit tests validate.

The contract is exact round-tripping:
$\text{dense (closed)} \to \text{packed} \to \text{dense (closed)}$ reproduces
the canonical matrix bit-for-bit, because unpacking plants the saved
constraints into an all-infinity matrix and re-runs the closure.  A debug mode
(\code{PTCTLMC\_CHECKINTERN}) asserts this on every intern.

The parallel-relevant part is not the byte count but \cref{lem:mingraph}: the
inclusion test can be decided \emph{on the compressed form} in $O(|G|)$.
Without it, answering ``is this successor already covered?'' requires an
$O(\dim^3)$ decompression of the stored zone --- which, in the two-phase
protocol, would be done while holding a lock.  On \code{arbiter17},
decompression was $36.3$\,s of work feeding $1.6$\,s of comparisons.

\subsection{Zone interning}

\FigIntern

Distinct zones are stored once and referenced by a $4$-byte index.  On
discrete-heavy symmetric models the sharing is extreme: \code{viking15} reaches
$2.77$\,M states drawn from about $120$ distinct zones, so the resident zone set
collapses from $\sim1.5$\,GB to a few kilobytes --- L2-resident.  On models
where nearly every zone differs (\code{fischer}, \code{arbiter}) interning costs
one hash per successor and loses nothing, because the copy had to be written
anyway.

The interner is sharded into $64$ independently locked partitions, and the
history of \emph{why} is instructive.  Interning sits on the per-successor hot
path, and the per-thread cache cannot absorb it: a symbolic run interns each
distinct zone exactly \emph{once}, so on \code{arbiter16} there were $589\,824$
distinct zones and a single cache hit in $589\,825$ calls --- essentially $100\%$
of calls reached the lock.  Hold time was tiny ($\sim0.2$ CPU-s) but the
\emph{acquisition rate} was not: threads waited about $20\times$ longer than the
lock was ever held, and the wait grew with thread count ($0.5 / 2.5 / 4.5$
CPU-s at $4/8/12$ threads), capping speedup at $\approx 5.5\times$.

Sharding fixed it, but only after one subtlety was resolved, which
\cref{prop:shard} formalises: the shard selector must not reuse the low hash
bits that also index a shard's hash table.  If it does, every key routed to
shard $s$ has the same low bits, so only $1/64$ of that shard's slots can ever
be a home position and the rest are reachable only by linear probing.  Measured
on \code{arbiter16}: $31.2$ probes per locked lookup versus $1.68$ unsharded ---
the sharding win, entirely eaten by probe chains.  Mixing and taking the
\emph{high} bits makes the two indices independent.

Above the shards sits a small per-thread direct-mapped cache
($2^{15}$ lines, $\approx 0.5$\,MB), which turns repeated interns into
lock-free hits.  A hit is verified by comparing the record \emph{bytes}, so a
$64$-bit hash collision cannot return a wrong index.  That verification is also
exactly why the cache must be dropped on GC ---
see \cref{prop:gc-cache}, which is a soundness obligation, not an optimisation.

\subsection{The frontier is not a list of states}

The frontier holds $(\dk, \mathit{zidx})$ pairs: $12$ bytes.  The location
vector and integer variables are decoded from $\dk$ on demand (the packing is a
lossless bijection), and the zone is decompressed from the interner into a
reused per-thread scratch at expansion time.  Storing full symbolic states
instead cost $\sim1.6$\,GB of zone copies on large \code{viking} and made the
frontier itself cache-hostile during the phase where all workers stream through
it.

\subsection{Active-clock reset}

Clocks that no process will read before resetting, and that the query does not
mention, are reset (\code{reset\_inactive}).  This collapses zones that differ
only in dead clocks, soundly --- it is the classical active-clock reduction of
Daws and Yovine \cite{DawsYovine96}.  The implementation is a bitmask fold: the
active set of a location vector is the OR of the per-location masks, so it costs
a handful of OR instructions plus a popcount loop instead of an
$\text{procs} \times \text{clocks}$ byte scan and a heap allocation --- on
\code{phil13} that ran $23$\,M times.
